\section{Supplementary results}\label{app:supp}

\paragraph{The native-pool round.} A second round (B1n--B4n) reran the same simulator with only the pool changed to 955,757 tokens. B2n and B3n were committed and pushed before any native E0b run had finished, 23 minutes after the native E0b jobs were submitted, and before any native E2 job existed, so B3n is blind in time and B2n is blind in that no result existed, not in submission order; their specification, like that of B1n and B4n, was written after the native curve was known (below), and B3n's cap of 64 is B4n's output. B2n held in part. Both modes predicted stability at all five rates where the engine was stable (10 of 10 seeds each), but the primary (mixed) mode put the all-stable edge at 6.5 and kept 10 of 10 runs stable there, where the real run was unstable, so its edge is too high by at least one step of its 0.5 grid in nominal multiples (the real 6.5 run's arrivals came to 6.86 $\lref$, above the real peak, and B4n puts the simulated peak 11\% above the real one); the secondary (separate) mode put the edge at 6.0, with 8 of 10 runs stable at 6.5, consistent with the real bracket. Because the real edge came in lower, the registered rates (2.5 and 4.0) lie below B3n's own test points, which are 6--16\% higher in multiples of $\lmiss$, and B3n was not simulated at them; 2.5 lies below every point it ran, toward the low end that it had flagged as its weak spot (9 of 10 seeds trapped at 2.85, the tenth recovering). The pass held at both registered rates. The levels are consistent with the lower rates, not a direct match: the trapped token rate rises with the rate in both the simulator and the engine, and the real 1.074--1.078 at 2.5 $\lref$ lies just below the 1.086--1.109 predicted at 2.65--2.85, while the real 1.106--1.124 at 4.0 overlaps the low end of the 1.119--1.140 predicted at 4.35--4.65. Outstanding backlogs of 549--594 and 1089--1130 requests lie inside the predicted 460--643 and 1067--1333 (they compare only roughly, as for B3), and good-arm hit fractions of 0.929--0.935 lie inside the predicted 0.928--0.936. B3n had also warned that at its higher point the trapped token rate sits close to the registered tolerance of 15\%, so that at higher rates the rule would stop calling a stuck run trapped; the real 1.106--1.124 at 4.0 $\lref$ is inside the tolerance, by 2.6 to 4.4 percentage points.

B1n and B4n are not blind. The native curve and $\lmiss$ had been measured and summarised in our project status notes (\fpath{notes/STATE.md}) before this round's specification was written, and the specification's finer grid near the cliff ($n=72$, 80, 88) may reflect that; only the LLM agent that ran the predictions was instructed not to read those data and, by its own record (not checked independently), did not. For completeness: B1n's $\lmiss$, 0.0720\sps{} and 11,239 tokens/s, is within 0.4\% of the real 0.0718 and 11,282. B4n put the peak at $n=72$, which the real grid does not contain, with $\mu$ 11\% above the real peak at 64--80; it put the 0.5-hit crossing at 94, against about 93, and the cap at 64, as registered. It flagged that cap as knife-edge ($\mu(64)$ only 0.35\% above the 5\% threshold, where a slightly lower value would have given 72) and the cliff cells $n=80$--96 as cells whose outcome differed across seeds; the real cliff lies somewhere between 80 and 96.

\paragraph{Stock LPM and its fallback.} In deep congestion stock LPM stops re-sorting its queue, so, on our reading of its source, Proposition~\ref{prop:A} applies to it (Appendix~\ref{app:proofs}); consistent with this, Stage~C's LPM run at $n=300$ had zero hits with a queue above 128, and round 5 of our earlier simulator, which modelled the 128-request fallback, found no LPM advantage in congestion. In the same earlier, uncalibrated simulator, pure LPM without the fallback has a congested service rate of about 1.17 times that simulator's own $\lmiss$ and recovered from 300 cold sessions only up to about 1.1 times it (\S\ref{sec:mech}); these are deep-congestion quantities and do not compare directly with recovery from a finite burst. We do not claim that LPM fails, and we have not isolated why it recovered. Post hoc, the timing matches the 128-request fallback: in the five runs that recovered, the engine's waiting queue (5\,s samples) never exceeded 128 after the burst (at most 116--126), and in each of the five that ended trapped the last hit fell in the 100\,s window in which the queue last stood at or below 128. This is a correlation in ten runs, not an isolation. Other candidates are partial prefix hits and LPM's own prefix matching, which moves the cached part of every waiting request's prefix to the recent end of the LRU order on each scheduling attempt; the simulator omits both. Under LPM the waiting-prefix refresh is skipped, because that matching already touches the waiting prefixes, so the two LPM configurations differ only while more than 128 requests wait and LPM has fallen back to FCFS. Seed 0 at 1.35 $\lref$ recovered with the refresh and was trapped without it; since the refresh made no difference under FCFS, single runs near LPM's threshold appear sensitive to small differences.

\paragraph{Trace-driven tool times.} TraceLab's sessions come from 52 pseudonymous users (release v0.0.2, a later snapshot than the one the paper describes)~\citep{zhu2026tracelab}. From TraceLab we take the per-turn tool time (mean 8.51\,s, median 0.237\,s, coefficient of variation 4.65) and, as an upper bound on it, the idle gap between requests (mean 14.52\,s). Both are capped at 600\,s, which affects 0.18\% of turns and 0.38\% of gaps but removes 38\% and 66\% of their raw means (13.6\,s and 42.9\,s); the per-turn series includes sub-agent runs and can include waits at permission prompts. Only TraceLab's tool times are used: its prompts average 172k tokens, more than the testbed's 49,152-token context limit. The mini-SWE-agent law comes from 500 published bash-only trajectories, the same 100 instances for each of five models (mean 1.40\,s, median 0.104\,s, coefficient of variation 4.29); 0.56\% of its steps hit the harness's 60\,s command limit, so the law is specific to that limit. The exponential control reproduced the record workload exactly and passed at both points. A secondary workload with mini-SWE-agent's mean turn and token counts (geometric turn counts) passed at 1.65 $\lmiss$ and failed at 2.35. There CAPPED (cap 12, the curve's peak) had left the trap, but at $\mu(12)-\lambda\approx0.022$\sps{} it needs about 9000\,s to drain 200 sessions, longer than the 7500\,s the run lasts after the burst; NOBURST was good in 8 of 10, and the other two collapsed without a burst. This workload's sessions outlive the declared measurement windows ($\lmiss$ is 0.0120\sps{} as declared, 0.0141 with a 20 times longer window, in which units the failed point is 1.99); a labelled post-hoc rerun with longer horizons gives, in those units, a narrower window, $[1,\,1.5]$, a pass at 1.35, and at 1.15 a failure of the slope rule (3 of 10 trapped, all 10 stuck below a hit fraction of 0.2).

\paragraph{Mechanisms in detail.} \emph{TTL pins.} Continuum's released preview pins for 2\,s, and only when the parsed tool's historical mean is at most 2\,s~\citep{li2025continuum}. We did not implement the paper's estimator. The preview's release rule is: when an allocation fails and nothing runs, release the pin that expires last; under it, at $\tau=20$\,s the system recovers from a cold start up to 4.0 $\lmiss$ and from 5 of 5 bursts. The rule we read in the paper's deadlock-prevention text is: when a request does not fit, unpin the latest-arriving programs until it fits, also while work runs; under it, $\tau=10$ and 20\,s recover from no cold start at 1.1 $\lmiss$ or above and from no burst. Our earlier rule, which waited for expiry while the GPU idled, gave 1.719 at $\tau=20$\,s. With plain FCFS order instead of Continuum's, under that earlier rule, the pins gave 0.29, 0.17 and 0.10 at $\tau=5$, 10 and 20\,s (3 seeds), with the GPU busy 33\%, 17\% and 8\% of the time: that simulator keeps an expired pin while the session's next request waits, which under FCFS lasts the rest of the queue wait. \emph{Adaptive caps.} With finished sessions' KV counted as usage, the 12 of 18 cells that stall do so at our surrogate's minimum window of 4, at 0.915--0.935 $\lmiss$; of five exploratory sensitivities, one (a 5\,s control interval instead of our surrogate's 10\,s) raised the rate by $0.184\pm0.167$, an unadjusted pointwise interval. Pause and resume of admitted agents is not modelled.
